\section*{Bloque B. Aplicaciones cuantitativas y optimización}

\noindent\textit{Preguntas 31 a 65. Requieren plantear y resolver problemas de producción, minimización de costos, maximización de beneficios y equilibrio competitivo con distintas formas funcionales, así como calcular productividades, elasticidades y economías de escala.}

\q{Con $Q=K^{0.5}L^{0.5}$, $K=16$ y $L=25$, la producción es:}
\begin{opts}
    \item 12
    \item 16
    \item 20
    \item 25
\end{opts}

\q{Con $Q=2K^{0.5}L^{0.5}$, $K=9$ y $L=4$, la producción es:}
\begin{opts}
    \item 12
    \item 24
    \item 6
    \item 18
\end{opts}

\q{Con $Q=K^{0.5}L^{0.5}$, $w=4$, $r=1$ y una meta de producción $Q=20$, el costo mínimo es:}
\begin{opts}
    \item 40
    \item 60
    \item 80
    \item 100
\end{opts}

\q{En el problema anterior, la relación óptima $K/L$ es:}
\begin{opts}
    \item 4
    \item $1/4$
    \item 2
    \item 1
\end{opts}

\q{Para $Q=K^{0.3}L^{0.6}$, la tecnología presenta rendimientos a escala:}
\begin{opts}
    \item crecientes (1.1).
    \item constantes (1.0).
    \item nulos (grado 0.0).
    \item decrecientes (0.9).
\end{opts}

\q{Para $Q=K^{0.4}L^{0.7}$, la tecnología presenta rendimientos a escala:}
\begin{opts}
    \item decrecientes.
    \item constantes.
    \item de 0.9.
    \item crecientes.
\end{opts}

\q{Con $Q=\min(2K,L)$, $w=2$, $r=3$ y $Q=40$, el costo mínimo es:}
\begin{opts}
    \item 100
    \item 120
    \item 140
    \item 160
\end{opts}

\q{Con $Q=\min(K,3L)$, $w=1$, $r=2$ y $Q=30$, el costo mínimo es:}
\begin{opts}
    \item 40
    \item 50
    \item 60
    \item 70
\end{opts}

\q{Con $Q=K+L$, $w=2$, $r=1$ y $Q=50$, el costo mínimo es:}
\begin{opts}
    \item 100
    \item 75
    \item 50
    \item 25
\end{opts}

\q{Con $Q=2K+L$, $w=3$, $r=2$ y $Q=40$, la cantidad óptima de capital es:}
\begin{opts}
    \item 10
    \item 15
    \item 20
    \item 40
\end{opts}

\q{En el corto plazo, $Q=10L^2-L^3$. El producto marginal del trabajo en $L=4$ es:}
\begin{opts}
    \item 32
    \item 24
    \item 16
    \item 8
\end{opts}

\q{Con la misma función $Q=10L^2-L^3$, el producto medio en $L=4$ es:}
\begin{opts}
    \item 16
    \item 24
    \item 32
    \item 20
\end{opts}

\q{Con $Q=10L^2-L^3$, el nivel de trabajo que maximiza el producto medio es:}
\begin{opts}
    \item 3
    \item 4
    \item 5
    \item 6
\end{opts}

\q{Si $C(Q)=100+4Q+Q^2$, el costo marginal en $Q=10$ es:}
\begin{opts}
    \item 4
    \item 14
    \item 20
    \item 24
\end{opts}

\q{Si $C(Q)=100+4Q+Q^2$, el costo medio en $Q=10$ es:}
\begin{opts}
    \item 14
    \item 20
    \item 22
    \item 24
\end{opts}

\q{Si $C(Q)=100+4Q+Q^2$, el costo variable medio en $Q=10$ es:}
\begin{opts}
    \item 4
    \item 14
    \item 24
    \item 10
\end{opts}

\q{Si $C(Q)=100+4Q+Q^2$, el nivel de producción que minimiza el costo medio es:}
\begin{opts}
    \item 5
    \item 20
    \item 14.14
    \item 10
\end{opts}

\q{Si $P=20$ y $C(Q)=100+4Q+Q^2$, la cantidad que maximiza el beneficio es:}
\begin{opts}
    \item 8
    \item 10
    \item 12
    \item 6
\end{opts}

\q{En el problema anterior, el beneficio económico en el óptimo es:}
\begin{opts}
    \item 36
    \item $-36$
    \item 0
    \item $-64$
\end{opts}

\q{Con $C(Q)=100+4Q+Q^2$, el precio de cierre o suspensión de la empresa en el corto plazo es:}
\begin{opts}
    \item 0
    \item 2
    \item 4
    \item 24
\end{opts}

\q{Suponga $C(Q)=100+4Q+Q^2$ para cada empresa y una demanda de mercado $Q=500-10P$. En el equilibrio de largo plazo, el número de empresas es:}
\begin{opts}
    \item 20
    \item 26
    \item 30
    \item 13
\end{opts}

\q{Si $P=20$, $CMg=4+2Q$ y la empresa produce $Q=8$, el excedente del productor es:}
\begin{opts}
    \item 32
    \item 64
    \item 96
    \item 128
\end{opts}

\q{Para $Q=(K^{0.5}+L^{0.5})^2$, la elasticidad de sustitución es:}
\begin{opts}
    \item 2
    \item 0.5
    \item 1
    \item infinita
\end{opts}

\q{Para $Q=(K^{-1}+L^{-1})^{-1}$, la elasticidad de sustitución es:}
\begin{opts}
    \item 0.0
    \item 1.0
    \item 2.0
    \item 0.5
\end{opts}

\q{Para $Q=K^{0.5}L^{0.5}$, con $w=4$ y $r=1$, la función de costo de largo plazo es:}
\begin{opts}
    \item $C=4Q$
    \item $C=2Q$
    \item $C=Q$
    \item $C=8Q$
\end{opts}

\q{Con $Q=K^{0.5}L^{0.5}$, $K=16$ fijo, $w=4$, $r=1$ y $Q=8$, el costo de corto plazo es:}
\begin{opts}
    \item 16
    \item 32
    \item 48
    \item 64
\end{opts}

\q{En el problema anterior, el costo marginal de corto plazo en $Q=8$ es:}
\begin{opts}
    \item 2
    \item 4
    \item 8
    \item 16
\end{opts}

\q{Si la función de costo de largo plazo es $C=4Q$, el costo marginal de largo plazo es:}
\begin{opts}
    \item 2
    \item 8
    \item 0
    \item 4
\end{opts}

\q{Suponga que el costo unitario es $C(Q)=100Q^{-0.2}$. Al duplicar la producción acumulada, el costo unitario se reduce aproximadamente en:}
\begin{opts}
    \item 5\%
    \item 10\%
    \item 13\%
    \item 20\%
\end{opts}

\q{Para $Q=K^{0.5}L^{0.5}$, la elasticidad del producto respecto del trabajo es:}
\begin{opts}
    \item 1
    \item 0.3
    \item 0.6
    \item 0.5
\end{opts}

\q{Si una tecnología tiene rendimientos crecientes a escala, la elasticidad del costo total respecto de la producción es:}
\begin{opts}
    \item mayor que 1.
    \item igual a 1.
    \item igual a 0.
    \item menor que 1.
\end{opts}

\q{Para $Q=K^{0.5}L^{0.5}$, la elasticidad de la demanda condicionada de trabajo respecto del precio del capital es:}
\begin{opts}
    \item $-0.5$.
    \item $0.5$.
    \item 1.
    \item $-1$.
\end{opts}

\q{Si $Q=AK^{0.5}L^{0.5}$ y la productividad total de los factores $A$ aumenta 10\%, manteniendo $K$ y $L$ constantes, la producción aumenta:}
\begin{opts}
    \item 5\%.
    \item 0\%.
    \item 10\%.
    \item 20\%.
\end{opts}

\q{Con $Q=K^{0.5}L^{0.5}$, $w=1$, $r=4$ y $Q=10$, el costo mínimo es:}
\begin{opts}
    \item 40
    \item 20
    \item 60
    \item 80
\end{opts}

\q{Con $Q=K^{0.5}L^{0.5}$, $K=16$ fijo, $w=4$, $r=1$ y $Q=20$, el costo de corto plazo es:}
\begin{opts}
    \item 100
    \item 116
    \item 125
    \item 80
\end{opts}
